% ==============================================================================
%  The attention map as a lattice operator — standalone report
%
%  Companion to glueball_report/glueball_spectroscopy.tex; self-contained.
%  Build: pdflatex (run twice for refs).
%  Figures are PNGs under results/<study>/; \graphicspath lists those dirs.
% ==============================================================================
\documentclass[11pt,a4paper]{article}

\usepackage[margin=2.6cm]{geometry}
\usepackage[T1]{fontenc}
\usepackage[utf8]{inputenc}
\usepackage{lmodern}
\usepackage{microtype}
\usepackage{amsmath,amssymb,mathtools}
\usepackage{graphicx}
\usepackage{booktabs}
\usepackage{caption}
\usepackage{xcolor}
\usepackage[colorlinks=true,linkcolor=blue!50!black,citecolor=blue!50!black,urlcolor=blue!50!black]{hyperref}

\graphicspath{{../}{../results/}{../results/sampler/}{../results/glueball/}{../results/attention/}{../results/wilson_regression/}}

% ---- shorthand ---------------------------------------------------------------
\newcommand{\Tr}{\operatorname{Tr}}
\newcommand{\Obar}{\bar{O}}
\newcommand{\Abar}{\bar{A}}
\newcommand{\meff}{m_{\mathrm{eff}}}
\newcommand{\latt}{\ell_{\mathrm{att}}}
\newcommand{\xiA}{\xi_A}
\newcommand{\That}{\hat{T}}

\title{The Attention Map as a Lattice Operator:\\
Correlation Length from a Gauge-Equivariant Network's Gaze}
\author{Francesco Passante}
\date{August 11, 2026}

\begin{document}
\maketitle

\begin{abstract}
A gauge-equivariant attention network computes its attention weights from a
\emph{gauge-invariant} score. This report takes that observation seriously and
follows it to its consequence: any reduction of an attention row to a single
number per site defines a local gauge-invariant scalar field, i.e.\ an ordinary
lattice operator, and if the network acts on one timeslice at a time then that
operator's zero-momentum connected correlator has an exact transfer-matrix
decomposition. \emph{The attention field therefore has a mass}, and that mass is
the theory's mass gap. We test this in three-dimensional $\mathbb{Z}_2$ gauge
theory near its second-order point, where the correlation length is tunable over
$\xi = 2.0 \to 5.6$ lattice spacings, using five ensembles ($24^2\times48$,
$N=2000$) and five per-timeslice variational operators trained on a
Rayleigh-quotient loss. On $1200$ configurations no checkpoint has seen, the
correlation length read out of the attention tracks the classical one at
\textbf{Pearson $0.9966$} over a factor $2.7$ of dynamic range, at $5$--$7\%$
precision on both sides, and --- the strongest part of the evidence --- it
reproduces a \emph{non-monotonic} excursion, $\xi(0.7585) > \xi(0.760)$,
measured on the same configurations by two independent operators. A cross-$\beta$
control matrix shows $\xiA$ is $11.9\times$ more sensitive to the evaluation
ensemble than to the training $\beta$, so nothing was memorised, and an explicit
configuration-scramble null returns no mass at all.

The result splits into two claims of different strength, and we are careful to
keep them apart. The \emph{structural} claim --- equivariant attention maps are
bona fide lattice operators, with a mass, an overlap, and a spectral
decomposition --- is established by $\xiA \approx \xi$, and holds already at
random initialisation. The \emph{learning} claim rests instead on the
ground-state overlap $A_0$: training raises it by $+0.11$ to $+0.27$
($5.4\sigma$ to $23.5\sigma$, correlated differences on shared configurations)
and raises the site-to-site fluctuation of the attention by one to one and a
half orders of magnitude. ``The network discovers the correlation length'' is
\emph{not} a supportable statement; ``training removes excited-state
contamination from the attention field'' is. We also record why two earlier
attempts at the same question failed --- both measured the mean attention
\emph{kernel}, a statistic bounded by the receptive radius $R$ and centred by
the ball geometry, rather than the attention \emph{field} --- and the four
analysis defects that had to be repaired before the numbers above could be
trusted, the largest being a generalized eigenproblem that was selecting a
near-null direction at every $\beta$.
\end{abstract}

\tableofcontents

% ==============================================================================
\section{The question, and why it is not a trivial one}
\label{sec:intro}
% ==============================================================================

The interpretability program of this thesis rests on one premise: in a
gauge-equivariant network, the attention map is not a diagnostic heat map but a
\emph{measurement}. The conference abstract commits to two investigations of
that premise --- whether attention localizes in topologically rich regions, and
whether the \textbf{attention range correlates with the physical correlation
length}. This report is about the second.

Stated naively the question is almost too easy to fail: train one operator per
inverse coupling $\beta$, measure how far each attention head looks, plot it
against the correlation length $\xi(\beta)$, and see whether the two rise
together. That programme was run twice and produced a negative both times
(Section~\ref{sec:failures}). The diagnosis --- that the failures were a property
of the \emph{statistic} and not of the physics --- is what makes the rest of the
report possible, so it is worth stating the trap early: a range statistic
computed from the \emph{lattice-averaged} attention kernel is bounded above by
the receptive radius $R$ and centred by the geometry of the offset ball, and it
discards precisely the site-to-site variation that distinguishes attention from
convolution.

The replacement statistic is the subject of Section~\ref{sec:claim}. Instead of
asking how \emph{wide} the attention is, we treat the attention itself as a
field on the lattice and ask what \emph{mass} it has. That question has a sharp,
falsifiable answer supplied by the transfer matrix, needs no reference to the
receptive field at all, and --- unlike the range --- does not exist for a
convolution.

% ==============================================================================
\section{Theoretical background}
\label{sec:theory}
% ==============================================================================

\subsection{Masses, correlators, and the mass gap}
\label{sec:masses}

Euclidean lattice field theory computes expectation values as statistical
averages over gauge-field configurations $U$ with weight $e^{-S[U]}$. For a
lattice with a periodic temporal extent $N_t$ and a transfer matrix
$\That = e^{-\hat{H}}$, the correlator of a zero-momentum operator
\begin{equation}
  \Obar(t) \;=\; \sum_{\vec{x}} O(t,\vec{x})
\end{equation}
admits the spectral decomposition
\begin{equation}
  C(\Delta) \;=\; \big\langle\, \delta\Obar(t+\Delta)\, \delta\Obar(t) \,\big\rangle
  \;=\; \sum_{n>0} \big| \langle 0 | \hat{O} | n \rangle \big|^2
        \Big( e^{-m_n \Delta} + e^{-m_n (N_t - \Delta)} \Big),
  \label{eq:spectral}
\end{equation}
where $\delta\Obar = \Obar - \langle\Obar\rangle$ (the vacuum subtraction is
mandatory for any operator with the vacuum quantum numbers, which is every
operator considered here) and the second exponential is the periodic image. At
large $\Delta$ the lightest state dominates, so
\begin{equation}
  C(\Delta) \;\xrightarrow{\ \Delta\ \text{large}\ }\;
  A\Big( e^{-m_0\Delta} + e^{-m_0(N_t-\Delta)} \Big),
  \qquad
  \xi \;=\; \frac{1}{m_0}.
\end{equation}
Two derived quantities are used throughout:
\begin{itemize}
  \item the \textbf{effective mass} $\meff(\Delta) = \log[C(\Delta)/C(\Delta+1)]$,
        which converges to $m_0$ from above --- every excited state in
        Eq.~\eqref{eq:spectral} adds a faster-decaying exponential, so a
        contaminated correlator decays too fast and reads $m$ \emph{high},
        hence $\xi$ \emph{low}. This sign matters in
        Section~\ref{sec:overshoot};
  \item the \textbf{ground-state overlap fraction}
        \begin{equation}
          A_0 \;=\; \frac{A\big(1 + e^{-m_0 N_t}\big)}{C(0)}
          \;\in\; [0,1],
          \label{eq:A0}
        \end{equation}
        the fraction of the operator's own variance carried by the ground
        state. This is Morningstar--Peardon's measure of operator quality, and
        the quantity a variational construction maximises. Two operators
        measuring the same physics must agree on $m_0$; they may differ
        arbitrarily on $A_0$, and that is where operator quality lives.
\end{itemize}

\subsection{The testbed: 3D $\mathbb{Z}_2$ gauge theory}
\label{sec:z2}

The measurement needs a correlation length that is (i) large compared with one
lattice spacing, so that a \emph{length} is resolvable at all on an integer
grid, and (ii) tunable over a real dynamic range, so that a correlation between
two lengths is more than a two-point coincidence. Three-dimensional
$\mathbb{Z}_2$ gauge theory delivers both: it is dual to the 3D Ising model and
has a second-order transition at $\beta_c \approx 0.7614$, where the mass gap
closes and $\xi = 1/m$ diverges as $\xi \sim (\beta_c-\beta)^{-\nu}$.

This is a deliberate change of system. The $SU(2)$ glueball ensembles on which
the rest of this thesis's spectroscopy is done have $\xi_s \le 1.1$ spatial
lattice spacings at every $\beta$ examined (Section~\ref{sec:failures}), which
makes the question unaskable there for reasons that have nothing to do with
attention. $\mathbb{Z}_2$ is also the project's declared debug testbed
($n_c = 1$, links $\pm1$) and is cheap. The price is that every conclusion below
is a statement about the \emph{method}, not about $SU(N)$.

Five ensembles were generated at
\begin{equation}
  \beta \in \{0.7450,\ 0.7520,\ 0.7560,\ 0.7585,\ 0.7600\},
  \qquad 24^2\times48, \quad N = 2000,
\end{equation}
spanning $\xi \approx 2.0 \to 5.6$. Their generation required an exact
$\mathbb{Z}_2$ heat-bath, $P(U=+1) = \sigma(2\beta s)$ with $s$ the staple sum:
the single-link Metropolis proposal is the flip $U \to -U$, whose acceptance
\emph{collapses} to $\approx 0.02$ as the configuration orders near $\beta_c$,
freezing the chain exactly where $\xi$ is largest. Even with the heat-bath,
critical slowing down ($\tau \sim \xi^z$, $z\approx2$) drives the integrated
autocorrelation time to $\tau_{\rm int} \approx 44$ sweeps, which fixes the
production $n_{\rm skip} = 200$.

\subsection{Gauge-equivariant attention, and why its weights are physical}
\label{sec:gemhsa}

The GELT block computes, for each query site $x$, a score against every site
$x+\Delta$ in the $L_1$ ball of Manhattan radius $R$. The key and value matrices
at $x+\Delta$ are first \emph{parallel-transported} to $x$ along the (averaged)
shortest lattice paths, and the score is the real trace
\begin{equation}
  s^{(\ell,h)}_{x\to x+\Delta}
  \;=\; \operatorname{Re}\Tr\!\big[\, Q^\dagger(x)\, \tilde{K}(x{+}\Delta) \,\big],
  \qquad
  \tilde{K}(x{+}\Delta) \;=\; T_\Delta(x)\, K(x{+}\Delta)\, T_\Delta^\dagger(x),
  \label{eq:score}
\end{equation}
followed by a softmax over $\Delta$ giving the attention weights
$\alpha^{(\ell,h)}_{x\to x+\Delta}$. Under a local gauge transformation
$\Omega(x)$, both $Q(x)$ and the transported $\tilde{K}$ conjugate by the
\emph{same} $\Omega(x)$, so the trace in Eq.~\eqref{eq:score} is
\textbf{invariant}, and therefore so is $\alpha$.

This is the entire foundation of what follows, and it is a property of the
architecture rather than of the training. It is also exactly what a
non-equivariant network does not have: there, the attention weights depend on
the gauge in which the configuration happens to be written, and no gauge-fixing
convention makes a gauge-dependent quantity into an observable.

\subsection{The claim: the attention field has a mass}
\label{sec:claim}

Because $\alpha_{x\to\cdot}$ is gauge invariant at every site, \emph{any}
reduction of the attention row to a scalar,
\begin{equation}
  A(x) \;=\; f\big(\alpha_{x\to\cdot}\big),
  \label{eq:reduction}
\end{equation}
is a gauge-invariant local scalar field built from the links --- a lattice
operator in the ordinary sense, with all the machinery of
Section~\ref{sec:masses} available to it. Three reductions are used here, chosen
to be qualitatively different rather than exhaustive:
\begin{equation}
  \underbrace{\alpha_{x\to x}}_{\texttt{self}},
  \qquad
  \underbrace{\sum_\Delta |\Delta|_1\, \alpha_{x\to x+\Delta}}_{\texttt{ell}},
  \qquad
  \underbrace{-\sum_\Delta \alpha_{x\to x+\Delta}\log\alpha_{x\to x+\Delta}}_{\texttt{ent}}.
\end{equation}
Note that \texttt{ell} is the \emph{local, per-site} version of the failed range
statistic: the trap of Section~\ref{sec:failures} is not the reduction, it is
the averaging that used to follow it.

The trained operators are \textbf{per-timeslice} by construction --- the
transfer-matrix bound that makes the Rayleigh loss a mass estimate requires
$\Obar(t)$ to depend on slice $t$ alone, so the network sees one timeslice at a
time and has no temporal receptive field. Consequently $A(t,\vec{x})$ depends on
slice $t$ alone, its zero-momentum sum $\Abar(t) = \sum_{\vec x} A(t,\vec x)$ is
a legitimate timeslice operator, and Eq.~\eqref{eq:spectral} applies verbatim:
\begin{equation}
  C_A(\Delta)
  \;=\; \big\langle \delta\Abar(t+\Delta)\, \delta\Abar(t) \big\rangle
  \;=\; \sum_{n>0} \big|\langle 0|\hat{A}|n\rangle\big|^2 e^{-m_n\Delta} + \dots
  \label{eq:CA}
\end{equation}
So the attention field has a mass $m_A$, and
\begin{equation}
  \boxed{\ \xiA \;=\; 1/m_A \;\overset{?}{=}\; \xi(\beta)\ }
\end{equation}
with \emph{no ceiling at $R$} and \emph{no geometric offset}, because
Eq.~\eqref{eq:CA} correlates a fluctuation --- zero-mean by construction --- and
not a moment of a kernel.

Three consequences, in increasing order of interest:
\begin{enumerate}
  \item \textbf{$\xiA$ tracks $\xi(\beta)$}: the length the failed statistic was
        supposed to find, recovered from the fluctuations of the attention
        rather than from its width.
  \item \textbf{This measurement does not exist for a convolution.} A fixed
        kernel has $\delta A \equiv 0$ identically, so $C_A \equiv 0$. The mere
        \emph{existence} of a signal is the content-dependence result, and it is
        quantitative: the site-level fluctuation-to-mean ratio $\delta A/A$ is
        its effect size.
  \item \textbf{$A_0$ of the attention field} says whether \emph{training}
        taught the routing to look at the physics, or whether the equivariant
        construction alone is responsible.
\end{enumerate}

\subsection{Why this is a measurement and not a tautology}
\label{sec:tautology}

Point~1 above is, on its own, close to a theorem: \emph{any} local
gauge-invariant scalar has a correlator that decays at the mass gap, a
random-init network's attention field included. Finding $\xiA \approx \xi$
therefore establishes point~2 --- that attention maps in an equivariant network
are physical fields --- but says nothing whatsoever about learning. The design
must separate the two, and does so with three arms:

\begin{center}
\small
\begin{tabular}{@{}l@{\quad}p{7.6cm}@{}}
\toprule
arm & what it isolates \\
\midrule
trained network, \textbf{matched} $\beta$ (diagonal)
  & the headline $\xiA$ vs $\xi$ \\[2pt]
trained network, \textbf{mismatched} $\beta$ (off-diagonal)
  & content-dependence vs memorisation: $\xiA$ must follow the
    \emph{evaluation} ensemble, not the training $\beta$ \\[2pt]
\textbf{random-init} network, all $\beta$
  & what equivariance gives \emph{before} learning: the $A_0$ baseline \\
\bottomrule
\end{tabular}
\end{center}

The off-diagonal arm also removes the confound that killed the previous attempt
(Section~\ref{sec:failures}): when one network per $\beta$ is read out on its
own ensemble, operator quality varies along the very axis being plotted. With
one fixed network evaluated across five ensembles, operator quality is a
property of the \emph{row} and the physics being probed is a property of the
\emph{column}.

% ==============================================================================
\section{Two prior negatives, and their shared cause}
\label{sec:failures}
% ==============================================================================

\subsection{$SU(2)$: the question is unaskable}

A classical GEVP $\beta$-scan on the $12^3\times24$, $\xi=3$ anisotropic
$SU(2)$ family (the same lattice on which this thesis's glueball result was
obtained) gives spatial correlation lengths
\begin{center}
\begin{tabular}{@{}lccccc@{}}
\toprule
$\beta$ & 2.1 & 2.3 & 2.4 & 2.5 & 2.7 \\
$\xi_s$ [$a_s$] & 0.57 & 0.83 & 1.00 & 1.10 & 1.00 \\
\bottomrule
\end{tabular}
\end{center}
$\xi_s$ never exceeds $1.1$ lattice spacings. Attention offsets are integers, so
the entire variation a range statistic would have to track lives inside the
first offset shell. The box is not the problem --- $L/\xi_s \approx 12$, the
regime in which finite-volume effects are \emph{absent} --- the spacing is: at
$\xi=3$ the spatial coupling $\beta_s = \beta/\xi \approx 0.8$ is strong
coupling, $a_s$ is coarse, and the whole glueball fits inside one grid cell.

\subsection{$\mathbb{Z}_2$: real dynamic range, and the statistic still fails}

Moving to 3D $\mathbb{Z}_2$ bought $\xi = 2.05 \to 5.28$, i.e.\ genuine range.
One variational operator per $\beta$ was trained, and the mean attention range
\begin{equation}
  \latt \;=\; \sum_\Delta |\Delta|_1\, \bar\alpha(\Delta),
  \qquad
  \bar\alpha = \text{attention averaged over sites and configurations}
\end{equation}
was read out. At $R=6$ it gave $4.25, 4.54$; at $R=12$ it gave
$8.63, 7.90, 8.13, 7.44, 8.21$ --- scatter, no trend in $\xi$.

The reason is that $\latt$ has a \textbf{uniform reference} that nobody had
quoted. For attention spread uniformly over the 2D $L_1$ ball,
\begin{equation}
  \ell_{\rm uniform} \;=\; \frac{\sum_d 4d\cdot d}{n_{\rm off}}
  \;=\; 4.28 \ \ (R=6), \qquad 8.31 \ \ (R=12),
\end{equation}
and \emph{every measurement landed on that number}. The statistic is bounded
above by $R$ \emph{and} centred by the ball geometry; only the bound had been
guarded against. In hindsight the $R=6$ ``trend'' ($4.25 \to 4.54$ about a
uniform $4.28$) was near-uniform noise and should never have been read as a
signal. A later reproducibility check settled it independently: retraining at
the same $R$ and the same $\beta$ with more data moved $\latt$ by $0.3$--$0.5$,
against an across-$\beta$ span of $3.95 \to 4.57$ for the whole scan. The
run-to-run variation of the statistic is the size of the entire effect it was
supposed to measure.

\subsection{The diagnosis}

Both failures share one cause. $\bar\alpha$ is an average \emph{over the
lattice}, so $\latt$ is a property of the learned \emph{kernel} --- a design
quantity of the same kind as a convolution's filter width --- and averaging is
exactly the operation that discards the site-to-site,
configuration-to-configuration variation that attention has and convolution does
not.

That the discarded part is large was already visible in the $R=12$ readout: the
per-head offset entropy relative to its uniform ceiling at $\beta=0.745$ reads
$0.87$, $0.92$, $0.93$, $0.97$, $\mathbf{0.42}$, $0.73$, $0.97$,
$\mathbf{0.72}$. Several heads are strongly structured --- while their $\latt$
sits at $8.3$ against a uniform $8.31$. The attention was doing something; the
radial first moment was blind to it. \textbf{Keeping the variation instead of
averaging it away is the whole content of Section~\ref{sec:claim}.}

% ==============================================================================
\section{What we did}
\label{sec:method}
% ==============================================================================

\subsection{Ensembles and trained operators}

Everything the measurement needs already existed when it was designed; no
sampling and no training were required, only forward passes.

\paragraph{Ensembles.} Five $\mathbb{Z}_2$ ensembles, $24^2\times48$, $N=2000$
configurations each, exact heat-bath, $n_{\rm skip} = 200$
(Section~\ref{sec:z2}).

\paragraph{Operators.} One GELT variational operator per $\beta$, trained
exactly as the $SU(2)$ glueball operator is: \texttt{GELT} with $4$ GEMHSA
layers, $2$ heads, $d_{\rm model}=16$, $d_{qkv}=8$, RoPE positional encoding,
receptive radius $R=6$, fed four cumulative APE-smearing levels
$(0,4,8,16)$ as input channels, and trained on the Rayleigh loss
$-\overline{C(\Delta)/C(0)}$ over $\Delta\in\{1,2\}$ whose converged value
\emph{is} (minus the exponential of) the mass. The operator is
per-timeslice --- 2D here, since 3D $\mathbb{Z}_2$ has two spatial directions
--- so the transfer-matrix bound of Section~\ref{sec:claim} holds.

Training used \texttt{configs[:800]} at every $\beta$. \textbf{All measurements
below use \texttt{configs[800:]}}, so the evaluation slice is unseen by every
checkpoint, the diagonal included.

\paragraph{Operator quality.} $R=6$ with $800$ training configurations was
chosen after an $R=12$ sweep produced systematically undertrained operators
(halving the data to pay for $313$ offsets left $70$ optimizer steps/epoch
instead of $140$). Each checkpoint is gated against the classical mass with
\emph{matched estimators} --- the network's own correlator re-fitted with the
classical scan's window and blocked jackknife, with both errors propagated ---
because the first version of the gate compared $\meff(\Delta{=}1)$ against a
cosh fit and thereby manufactured an apparent ``undertraining that grows with
$\xi$'' out of a purely one-sided estimator bias. All five $R=6$ checkpoints
pass the corrected gate.

\subsection{The measurement pipeline}
\label{sec:pipeline}

For each (network, ensemble) pair, and for each of the $4\ \text{layers} \times
3\ \text{reductions} \times 2\ \text{heads} = 24$ attention channels:

\begin{enumerate}
  \item run the forward pass on the unseen configurations and stash
        \texttt{\_last\_alpha} per layer (under \texttt{no\_grad});
  \item reduce each attention row to a scalar per site via
        Eq.~\eqref{eq:reduction}, and accumulate the site-level first and second
        moments so the fluctuation-to-mean ratio $\delta A/A$ can be quoted;
  \item form the zero-momentum series $\Abar(t)$ per configuration;
  \item project the multi-channel basis onto its GEVP ground-state vector
        ($t_0=1$, $t_d=2$), with the safeguards of
        Section~\ref{sec:defects};
  \item fit the projected connected correlator to the periodic single-state
        form over $\Delta\in[2,8]$;
  \item take errors from a \textbf{blocked} jackknife over configurations with
        block $20$.
\end{enumerate}

Conventions in steps 4--6 are copied verbatim from the classical
$\beta$-scan so the two masses are directly comparable, and the classical
smeared-plaquette basis (levels $0,4,8,16$) is \textbf{re-measured on the same
configurations}, so $\xiA - \xi_{\rm class}$ can be quoted as a correlated
difference rather than as two independent numbers.

\paragraph{Reference arms.} Alongside the attention channels we measure the
network's own output operator $\Obar(t)$ (a check that the checkpoint loaded
correctly: its mass must reproduce the recorded one) and the classical basis.

\paragraph{Nulls.} Two, and the distinction between them turned out to matter
(Section~\ref{sec:defects}):
\begin{itemize}
  \item \textbf{config-scramble} --- each timeslice is drawn from an
        independently permuted configuration. This destroys the temporal
        correlation \emph{and} the per-configuration zero mode, so $C(\Delta>0)$
        has expectation zero and there is nothing for a mass to be fitted to.
        This is the null proper.
  \item \textbf{time-shuffle} --- one independent permutation of the timeslices
        per configuration. This is \emph{not} a null (it preserves each
        configuration's mean over $t$); it is retained as a
        \textbf{zero-mode consistency check}, whose plateau must equal
        $(2\xi-1)/N_t$.
\end{itemize}

\paragraph{Cost.} Two runs. The \emph{diagonal} run (matched network + random
control, $N_{\rm eval} = 1200$ configurations per ensemble) buys precision where
precision is the point. The \emph{cross} run (all five networks + random on all
five ensembles, $N_{\rm eval} = 600$) buys the control matrix, where only the
row/column contrast matters. About four minutes per ensemble per network.

% ==============================================================================
\section{Results}
\label{sec:results}
% ==============================================================================

\subsection{The headline: the attention field carries the mass gap}
\label{sec:headline}

Table~\ref{tab:headline} is the diagonal run: every trained network on its own
ensemble, on the $1200$ configurations it never saw, against the classical basis
measured on those same configurations.

\begin{table}[htbp]
\centering
\caption{Correlation lengths from the attention field, from the classical
smeared basis, and from a randomly initialised network of identical
architecture. $1200$ unseen configurations per ensemble; errors are blocked
jackknife (block $20$). ``channel'' is the attention channel the variational
self-check selected (layer, head, reduction). $\xi_{\rm out}$ is the network's
own output operator, shown for reference.}
\label{tab:headline}
\begin{tabular}{@{}lccccc@{}}
\toprule
$\beta$ & $\xi$ classical & $\xiA$ \textbf{trained} & channel
        & $\xiA$ random init & $\xi_{\rm out}$ \\
\midrule
0.7450 & 2.04 (14) & 2.24 (12) & \texttt{L4h0:self} & 2.08 (13) & 2.33 \\
0.7520 & 2.34 (16) & 2.76 (13) & \texttt{L4h1:self} & 2.46 (16) & 2.92 \\
0.7560 & 4.08 (27) & 4.48 (30) & \texttt{L4h0:ell}  & 4.08 (26) & 4.63 \\
0.7585 & \textbf{5.59 (34)} & \textbf{6.64 (32)} & \texttt{L4h0:ent} & 5.47 (31) & 6.65 \\
0.7600 & \textbf{4.71 (34)} & \textbf{5.60 (37)} & \texttt{L4h1:ent} & 4.80 (34) & 5.81 \\
\bottomrule
\end{tabular}
\end{table}

\begin{equation}
  \textrm{Pearson}\big(\xiA,\ \xi_{\rm classical}\big) \;=\; \mathbf{0.9966},
  \qquad \textrm{slope } 1.21,
\end{equation}
over a factor $2.74$ of classical dynamic range, at $5$--$7\%$ precision on both
sides. The claim of Section~\ref{sec:claim} holds: the attention map is a
lattice operator, and its correlator decays with the mass gap.

\paragraph{The correlation survives a non-monotonic excursion.} This is the
strongest single piece of evidence in the report and it deserves to be read
carefully. $\xi(0.7585) > \xi(0.760)$ --- i.e.\ the point \emph{further} from
$\beta_c$ has the \emph{larger} correlation length --- in the classical
measurement \emph{and} in the attention, on the same configurations, from
independent operators. Two quantities that both rise with $\beta$ correlate for
trivial reasons; two that both turn around at the same place do not. (An earlier
$N=2000$ scan had the opposite ordering at these two points; those were its two
worst-determined points, measured with the defective GEVP of
Section~\ref{sec:defects}, and they are superseded.) The excursion is a property
of the \emph{ensembles}, not of the attention: both points sit at
$\xi \gtrsim L/4 = 6$, the scan's own finite-volume guard --- see
Section~\ref{sec:limits}.

\paragraph{Correlator profiles.} At $\beta=0.760$, the point closest to
$\beta_c$, $C(\Delta)/C(0)$ reads
\begin{center}
\begin{tabular}{@{}lccccccc@{}}
\toprule
 & $\Delta=0$ & 1 & 2 & 3 & 4 & 5 & 6 \\
\midrule
classical smeared basis  & $+1.000$ & $+0.404$ & $+0.272$ & $+0.197$ & $+0.163$ & $+0.134$ & $+0.109$ \\
\textbf{trained attention} & $+1.000$ & $+0.619$ & $+0.470$ & $+0.373$ & $+0.312$ & $+0.260$ & $+0.224$ \\
random attention         & $+1.000$ & $+0.466$ & $+0.310$ & $+0.226$ & $+0.186$ & $+0.155$ & $+0.128$ \\
config-scramble null     & $+1.000$ & $+0.005$ & $+0.007$ & $+0.009$ & $+0.010$ & $-0.007$ & $-0.000$ \\
\bottomrule
\end{tabular}
\end{center}
The trained attention field is a visibly cleaner exponential than either the
classical basis or the untrained network, and the null is flat at the $10^{-2}$
level with no fittable mass.

\subsection{Training separates from initialisation at every $\beta$}
\label{sec:learning}

Table~\ref{tab:learning} is the learning claim. $\Delta A_0$ is a blocked
jackknife of the \emph{difference} on shared configurations: both arms see
identical configurations, so the naive $\sqrt{\sigma^2+\sigma^2}$ throws away
the cancellation.

\begin{table}[htbp]
\centering
\caption{Ground-state overlap of the attention field, trained vs randomly
initialised, and the site-level fluctuation-to-mean ratio $\delta A/A$ (median
over the $24$ channels). $\Delta A_0$ and $\Delta\xi$ are correlated
differences on shared configurations.}
\label{tab:learning}
\footnotesize
\setlength{\tabcolsep}{4pt}
\begin{tabular}{@{}lcccccc@{}}
\toprule
$\beta$ & $A_0$ trained & $A_0$ random & $A_0$ classical
        & $\delta A/A$ (tr / rnd) & $\Delta A_0$ (correlated) & $\Delta\xi$ \\
\midrule
0.7450 & 0.848 (41) & 0.704 (42) & 0.578 (39) & 0.091 / 0.0067 ($\times14$) & $+0.144\pm0.017$ (\textbf{8.4$\sigma$})  & $+0.17\pm0.04$ \\
0.7520 & 0.782 (28) & 0.671 (36) & 0.576 (37) & 0.140 / 0.0060 ($\times23$) & $+0.110\pm0.020$ (\textbf{5.4$\sigma$})  & $+0.31\pm0.07$ \\
0.7560 & 0.738 (17) & 0.534 (16) & 0.449 (15) & 0.246 / 0.0054 ($\times46$) & $+0.204\pm0.012$ (\textbf{16.9$\sigma$}) & $+0.41\pm0.13$ \\
0.7585 & 0.753 (11) & 0.486 (14) & 0.413 (13) & 0.113 / 0.0049 ($\times23$) & $+0.267\pm0.011$ (\textbf{23.5$\sigma$}) & $+1.17\pm0.19$ \\
0.7600 & 0.654 (14) & 0.449 (15) & 0.397 (14) & 0.130 / 0.0046 ($\times28$) & $+0.205\pm0.012$ (\textbf{16.9$\sigma$}) & $+0.80\pm0.19$ \\
\bottomrule
\end{tabular}
\end{table}

Three things to read off it. First, the trained attention field is a
\emph{better operator than the classical smeared basis} at every $\beta$, by
$0.20$ to $0.34$ in overlap fraction --- and the classical basis degrades
towards $\beta_c$ ($0.578 \to 0.397$) faster than the attention does
($0.848 \to 0.654$). Second, the random network's site-to-site fluctuation is
$\approx 0.005$ essentially independently of $\beta$, and training raises it by
one to one and a half orders of magnitude; for a fixed convolutional kernel the
same quantity is identically zero. Third, $\Delta\xi > 0$ everywhere: the
trained attention field is also \emph{less} contaminated by excited states than
the untrained one, not merely better normalised.

\paragraph{Autocorrelation.} $\tau_{\rm int}$ of the selected channel along the
chain reads $0.63$, $1.13$, $2.65$, $\mathbf{8.35}$, $5.85$ across the five
$\beta$. The jackknife block of $20$ clears $2\tau_{\rm int}$ at every $\beta$,
but only just at $\beta=0.7585$.

\paragraph{The zero-mode check.} The time-shuffled arm's plateau matches its
prediction $(2\xi-1)/N_t$ at every $\beta$: measured $0.062$, $0.079$, $0.138$,
$0.231$, $0.159$ against predicted $0.073$, $0.094$, $0.166$, $0.256$, $0.212$,
tracking at Pearson $0.98$ \emph{including} the non-monotonic excursion. This is
a second, independent confirmation of $\xi$ from the same data.

\subsection{The cross-$\beta$ control matrix: no memorisation}
\label{sec:cross}

Five trained networks plus the random arm, each evaluated on all five ensembles,
$600$ unseen configurations per cell. One cell (\texttt{train@0.756} evaluated on
$\beta=0.752$) returned $A_0 = 5.76$ with $\xi = 0.39$ --- a fit that ran to its
grid edge --- and is masked throughout; the guard now reports such cells as
unresolved rather than as numbers.

\begin{table}[htbp]
\centering
\caption{$\xiA$ from the cross-evaluation matrix. Rows are the $\beta$ at which
the network was trained; columns are the ensemble it is evaluated on. The
masked cell is marked ---.}
\label{tab:cross}
\begin{tabular}{@{}lccccc@{}}
\toprule
network & \multicolumn{5}{c}{evaluation ensemble $\beta$} \\
\cmidrule(l){2-6}
trained at & 0.7450 & 0.7520 & 0.7560 & 0.7585 & 0.7600 \\
\midrule
0.7450 & 2.24 & 2.73 & 4.53 & 5.73 & 4.85 \\
0.7520 & 2.33 & 2.77 & 4.64 & 5.74 & 4.92 \\
0.7560 & 2.36 & --- & 4.71 & 6.06 & 5.04 \\
0.7585 & 2.31 & 2.80 & 4.68 & 6.50 & 5.25 \\
0.7600 & 2.22 & 2.69 & 4.68 & 6.15 & 5.13 \\
\midrule
random init & 2.05 & 2.47 & 4.32 & 5.42 & 4.77 \\
\midrule
\emph{classical} & 2.05 & 2.36 & 4.21 & 5.48 & 4.73 \\
\bottomrule
\end{tabular}
\end{table}

The matrix reads across, not down:
\begin{equation}
  \frac{\text{spread within a row (across evaluation ensembles)}}
       {\text{spread within a column (across training $\beta$)}}
  \;=\; \frac{1.405}{0.118} \;=\; \mathbf{11.9},
\end{equation}
and the column means correlate with the classical $\xi$ at \textbf{Pearson
$0.9983$}. The attention reads the configuration in front of it, and a fixed
kernel has no analogue of this measurement at all.

\paragraph{The transfer test.} If a network had learned its own ensemble rather
than physics, the diagonal would dominate its column. It does not:

\begin{center}
\begin{tabular}{@{}lcccc@{}}
\toprule
$\beta$ & $A_0$ diagonal & $A_0$ off-diagonal & advantage & $A_0$ random \\
\midrule
0.7450 & 0.855 & 0.770 & $+0.085$ & 0.723 \\
0.7520 & 0.774 & 0.815 & $\mathbf{-0.041}$ & 0.656 \\
0.7560 & 0.697 & 0.687 & $+0.010$ & 0.484 \\
0.7585 & 0.764 & 0.696 & $+0.069$ & 0.495 \\
0.7600 & 0.660 & 0.640 & $+0.020$ & 0.423 \\
\bottomrule
\end{tabular}
\end{center}

The mean diagonal advantage is $+0.029$ and is \emph{inconsistent in sign},
against a mean trained-minus-random of $+0.167$: the learning effect is
$\mathbf{6\times}$ the $\beta$-matching effect. Every trained network beats the
random arm on every ensemble --- the worst trained network on each ensemble
scores $0.747 / 0.774 / 0.636 / 0.659 / 0.571$ against the random arm's
$0.723 / 0.656 / 0.484 / 0.495 / 0.423$.

What the matrix \emph{does} separate is network quality: ensemble-averaged $A_0$
per row is $0.734 / 0.677 / 0.702 / \mathbf{0.775} / 0.726$ against $0.556$ for
random, so \texttt{train@0.7585} is the best operator on three of five
ensembles, including ones it never saw --- a globally better operator, not a
$\beta$-specialist. \textbf{The learning claim is therefore about transferable
physics}: the routing did not memorise a correlation length, it learned
something that works at every $\xi$ in the scan.

\paragraph{The null passes.} With the config-scramble, $26$ of the $30$ cells do
not fit at all --- there is nothing to fit --- and the four that do return
$|A_0| \le 0.0066$.

\begin{figure}[htbp]
\centering
\includegraphics[width=\textwidth]{z2_attention_correlator_R6.png}
\caption{The four panels of the cross-evaluation run. \textbf{Top left}: $\xiA$
against the classical $\xi$ measured on the same configurations, for the matched
trained network and for random initialisation; the dotted line is $y=x$.
\textbf{Top right}: the cross-evaluation matrix of Table~\ref{tab:cross} ---
$\xiA$ follows the column (the physics), not the row (the training $\beta$).
\textbf{Bottom left}: both lengths against $\beta_c-\beta$ with an effective
exponent; this is \emph{not} a measurement of $\nu$ (Section~\ref{sec:limits}).
\textbf{Bottom right}: ground-state overlap $A_0$ for the trained attention, the
random attention, and the classical basis.}
\label{fig:main}
\end{figure}

% ==============================================================================
\section{Analysis}
\label{sec:analysis}
% ==============================================================================

\subsection{What is established, and what is not}
\label{sec:split}

The random-init arm \emph{also} tracks $\xi$, at Pearson $0.9989$ --- if
anything more tightly than the trained arm. This is exactly what
Section~\ref{sec:tautology} predicted, and it is not a disappointment; it is the
correct partition of two claims of very different strength:

\begin{itemize}
  \item \textbf{Structural claim (established).} Attention maps in a
        gauge-\emph{equivariant} network are bona fide lattice operators, with a
        mass, an overlap, and a spectral decomposition. This is established by
        $\xiA \approx \xi$, and it is true already at initialisation. It is the
        premise of the whole interpretability program made quantitative, and it
        is a property of the \emph{architecture}.
  \item \textbf{Learning claim (established, but by a different statistic).}
        Training makes those operators \emph{good}. This is established by
        $A_0$ ($+0.11$ to $+0.27$, up to $23.5\sigma$), by $\delta A/A$
        ($\times14$ to $\times46$), and by the transfer test of
        Section~\ref{sec:cross} --- \emph{not} by $\xiA$.
\end{itemize}

\begin{quote}
``The network discovers the correlation length'' is \textbf{not} a supportable
claim and must not be made.
\end{quote}

\paragraph{A known confound on $\delta A/A$.} $\alpha$ is a softmax output, so a
network with larger score magnitudes has a sharper --- hence more variable ---
attention regardless of what it learned. $\delta A/A$ is therefore
\emph{descriptive}; the load-bearing learning statistic is $A_0$, which a
sharp-but-untrained attention would not raise (it would build a UV-dominated
field with \emph{poor} ground-state overlap). Closing this properly needs an
entropy-matched random arm, which has not been run.

\subsection{Why the trained $\xiA$ overshoots, and why ``closer to classical'' is the wrong target}
\label{sec:overshoot}

The natural inference from Section~\ref{sec:learning} is: the trained attention
has the better ground-state overlap, so its $\xiA$ should sit \emph{closer} to
the classical $\xi$. The data say the opposite, and the wrong reading here would
invert the conclusion.

\begin{center}
\begin{tabular}{@{}lccccc@{}}
\toprule
$\beta$ & $\xi$ classical & $\xiA$ trained & $\xiA$ random
        & $|{\rm trained}-{\rm class}|$ & $|{\rm random}-{\rm class}|$ \\
\midrule
0.7450 & 2.04 & 2.24 & 2.08 & 0.20 & \textbf{0.04} \\
0.7520 & 2.34 & 2.76 & 2.46 & 0.42 & \textbf{0.12} \\
0.7560 & 4.08 & 4.48 & 4.08 & 0.40 & \textbf{0.00} \\
0.7585 & 5.59 & 6.64 & 5.47 & 1.05 & \textbf{0.12} \\
0.7600 & 4.71 & 5.60 & 4.80 & 0.89 & \textbf{0.09} \\
\bottomrule
\end{tabular}
\end{center}

The \emph{random} arm is the one that agrees with the classical operator, at
every $\beta$ (fitted slope $0.96$ against the trained arm's $1.21$). This is
not a contradiction, because \textbf{the classical operator is not ground
truth}. It is one more operator with its own excited-state contamination, and
contamination has a \emph{signed} effect: by Eq.~\eqref{eq:spectral}, excited
states add faster-decaying exponentials, so a contaminated correlator decays
faster than the true gap and reads $\xi$ \textbf{low}. $A_0$ measures how much
of an operator's own $C(\Delta)$ is the ground state, so a higher $A_0$ means
less of that downward bias, and the operator with the highest $A_0$ should read
the \emph{largest} $\xi$. It does. The observed ordering
\begin{equation}
  \xi_{\rm trained} \;>\; \xi_{\rm random} \;\approx\; \xi_{\rm classical}
\end{equation}
is precisely the ordering the $A_0$ column predicts, once one stops treating
``classical'' as the target. The correlator profiles in
Section~\ref{sec:headline} say the same thing directly: at $\beta=0.760$ the
classical basis has fallen to $0.404$ by $\Delta=1$ and the random attention to
$0.466$, while the trained attention still holds $0.619$.

So the natural reading of the whole result is that the trained operator's larger
$\xi$ is the \emph{candidate} for the value closest to the true gap, and that
\textbf{learning $=$ removing excited-state contamination} is the physically
honest gloss of the $A_0$ result. That gloss is stronger than the correlation
framing, because it says what the training bought in operator terms --- which is
exactly what a variational-operator audience asks for.

\paragraph{Stated as a candidate, not as a result.} There is no independent
ground truth for $\xi$ at these $\beta$: the older $N=2000$ scan is superseded
by the GEVP defect below, and both top points sit at $\xi \gtrsim L/4$ where
finite volume bites. A $21\%$ systematic in the attention correlator fit would
look identical. Distinguishing ``de-contamination'' from ``systematic
overshoot'' needs a reference that is not itself contaminated. Consequently:
\begin{itemize}
  \item do \textbf{not} quote $|\xiA - \xi_{\rm classical}|$ as an accuracy
        metric, and do not present the random arm's smaller residual as the
        random arm being \emph{better} --- agreement with a contaminated
        reference is not accuracy;
  \item the load-bearing statements remain the \emph{correlation} (Pearson
        $0.9966$, through the non-monotonic excursion) for the structural claim
        and $A_0$ for the learning claim. The slope $1.21$ is a separate, open
        quantity.
\end{itemize}

\subsection{Four defects, and how each was resolved}
\label{sec:defects}

None of the numbers above survived the first version of the analysis. The
repairs are recorded because two of them are of general interest.

\paragraph{1. The GEVP was selecting a near-null direction --- at every $\beta$.}
The first suspicion was our own basis pruning; fixing it moved the classical
masses by $4\times10^{-15}$. That non-result \emph{is} the diagnosis: $m$, $A_0$
and $C(\Delta)/C(0)$ are all invariant under rescaling the projected operator,
so identical values mean the same direction was selected either way. The real
cause was an eigenvalue floor of $10^{-12}$ relative to the largest in the
whitening step --- twelve orders of magnitude, i.e.\ no regularisation at all. A
nested-smearing $C(t_0)$ has three or four meaningful orders, and a small signal
divided by a floored \emph{noise} eigenvalue produces a spuriously large
generalized eigenvalue pointing into noise: a near-cancelling combination whose
$C(0)$ is almost pure noise while $C(\Delta>0)$ still carries the physical
decay. The symptom was $C(1)/C(0) = 0.039 / 0.020 / 0.007$ at the three largest
$\beta$ with $36\%$ / $87\%$ / unresolved errors --- \emph{correct masses,
useless error bars}.

Fixed with a floor of $10^{-4}$ plus a \textbf{variational self-check}: $v_0$
maximises the Rayleigh quotient over the span of the basis, so the projection
can never interpolate worse than any single member; when it does, fall back to
that member. The decision is taken once on the full sample and held fixed across
jackknife replicas (a jackknife whose estimator changes between replicas does
not estimate anything). All five $\beta$ then resolve at $6$--$7\%$ with
$C(1)/C(0) = 0.39$--$0.44$ uniformly --- and the \textbf{fallback fires at every
$\beta$}, on both the classical and the attention bases.

\paragraph{2. The multi-channel attention GEVP is unstable.} Three NaN cells and
one absurd one ($\xi = 4.85 \pm 89$) out of ten. Moot after the self-check: it
falls back to the single best channel at every $\beta$, so the multi-channel and
best-single arms coincide and the tables quote one number. Informative in
itself --- the $24$ attention channels are redundant enough that a variational
combination never helps.

\paragraph{3. $\Delta A_0$ errors were uncorrelated.} Fixed with a blocked
jackknife of the difference on shared configurations. The two weakest points
went from $2.4\sigma$ to $8.4\sigma$ and $5.4\sigma$; nothing else moved.

\paragraph{4. A time-shuffle is not a null.} The null's $A_0$ rose from
$0.005$--$0.014$ to $0.043$--$0.230$ once the GEVP fallback landed. The first
suspicion was selection bias (the null re-selecting its own best-of-six; that
bias is real, and measures $7.4\times$ on pure noise). Removing it changed the
null by $0.000$. The correct explanation is structural: a permutation of
timeslices preserves each configuration's mean over $t$, while the connected
correlator subtracts only the \emph{global} mean, so the residual retains
\begin{equation}
  \operatorname{Var}_{\rm config}(\Obar_c) \;=\; C(0)\cdot\frac{2\xi-1}{N_t}
\end{equation}
as a \emph{flat pedestal at every $\Delta$} --- which is exactly what the
profiles show ($+1.000$, $+0.236$, $+0.229$, $+0.231$, $+0.234$, \ldots at
$\beta=0.7585$: no decay at all). The ``null'' was measuring the finite-$N_t$
zero mode --- physics, not noise --- and could never have gone to zero. Fixed
with the config-scramble of Section~\ref{sec:pipeline}, and validated on a
synthetic operator of known mass ($\xi = 5.56$ injected): the real arm returns
$\xi = 5.65$ with $A_0 = 1.000$, the time-shuffle plateau lands at $0.213$
against a predicted $0.215$, and the config-scramble returns $\pm0.001$ with no
fittable mass. The time-shuffle is retained as the zero-mode consistency check
of Section~\ref{sec:learning}.

\subsection{Limitations, stated up front}
\label{sec:limits}

\begin{enumerate}
  \item \textbf{One theory, one volume, one lattice family.} 3D $\mathbb{Z}_2$
        at $24^2\times48$. This is a statement about the \emph{method}, not
        about $SU(N)$. The natural next test is the $SU(2)$ glueball operator,
        where the same readout machinery already exists.
  \item \textbf{Finite volume at the top two $\beta$ is the open physics
        question.} $\xi = 5.6$ and $4.7$ against $L/4 = 6$ puts both points at
        or past the scan's own guard, and the non-monotonic excursion is what a
        volume-squeezed pair looks like. The test is $L = 24 \to 32$ --- but
        GELT is built per-$L$, so that means new ensembles \emph{and} new
        checkpoints at all five $\beta$ ($\approx2.4\times$ per sweep, tens of
        hours), not one point. Worth it only if the excursion is to be reported
        as physics rather than flagged as unresolved. It is flagged as
        unresolved here.
  \item \textbf{The effective exponent is not $\nu$.} Fitting
        $\xi \sim (\beta_c-\beta)^{-\nu}$ to these five points returns
        $0.35$--$0.48$ depending on the arm and the configuration budget
        ($0.43$ classical and $0.48$ attention on the diagonal run, $0.37$ and
        $0.35$ on the cross run of Figure~\ref{fig:main}), against the 3D Ising
        $\nu = 0.63$: $\xi \le 5.6$ at $L=24$ is not the asymptotic scaling
        regime, and the non-monotonic excursion sits inside the fit. It may be
        quoted only as a \emph{consistency} check between the attention and the
        classical scan on the same points --- and with the excursion now
        confirmed on both operators, it should probably not be quoted at all
        until the volume question is settled.
  \item \textbf{Operator quality is uneven across the five checkpoints.} The
        cross-evaluation design makes this a nuisance parameter rather than a
        confound (quality is a property of the row, physics of the column), but
        the diagonal alone would still inherit it.
  \item \textbf{The $\delta A/A$ confound} of Section~\ref{sec:split}: no
        entropy-matched random arm has been run.
  \item \textbf{One masked cell} in the cross matrix (Section~\ref{sec:cross}),
        excluded by an automatic guard rather than by hand, but excluded
        nonetheless.
\end{enumerate}

% ==============================================================================
\section{Conclusions}
\label{sec:conclusions}
% ==============================================================================

The question the conference abstract poses --- does the attention range track
the physical correlation length? --- was closed negative twice, and both
negatives were about the statistic rather than the physics. $\latt$, the first
radial moment of the \emph{lattice-averaged} attention, is bounded above by the
receptive radius and centred by the geometry of the offset ball; it is a
property of the learned kernel, of the same kind as a convolution's filter
width, and the averaging that defines it discards exactly what distinguishes
attention from convolution. \textbf{Any range statistic must be quoted against
its uniform reference}, and this one never survives that comparison.

Keeping the site-to-site variation instead turns the question into a
well-posed one. Because the equivariant score is gauge invariant, any reduction
of an attention row is a local gauge-invariant scalar field; because the trained
operator is per-timeslice, that field's zero-momentum correlator has an exact
transfer-matrix decomposition. The attention field therefore \emph{has a mass},
and it is the theory's mass gap. Measured on $1200$ unseen configurations per
ensemble across five 3D $\mathbb{Z}_2$ ensembles:

\begin{itemize}
  \item $\xiA$ tracks the classical $\xi$ at \textbf{Pearson $0.9966$} over
        $\times2.7$ of dynamic range at $5$--$7\%$ precision on both sides, and
        reproduces a non-monotonic excursion in $\beta$ that two independent
        operators see on the same configurations;
  \item $\xiA$ is $\mathbf{11.9\times}$ more sensitive to the evaluation
        ensemble than to the training $\beta$, with column means correlating
        with the classical $\xi$ at Pearson $0.9983$, so the networks read the
        configuration in front of them rather than recalling their own;
  \item a config-scramble null returns no mass in $26$ of $30$ cells and
        $|A_0| \le 0.0066$ in the rest;
  \item training raises the ground-state overlap of the attention field by
        $+0.11$ to $+0.27$ ($5.4\sigma$ to $23.5\sigma$) over random
        initialisation, and raises its site-level fluctuation by one to one and
        a half orders of magnitude --- while the same fluctuation is identically
        zero for a fixed convolutional kernel, so this measurement has no
        analogue there at all.
\end{itemize}

The honest partition of these results is that the \emph{structural} claim
--- gauge-equivariant attention maps are lattice operators, with a mass, an
overlap, and a spectral decomposition --- is established, and holds already at
initialisation; while the \emph{learning} claim is established by $A_0$ rather
than by $\xiA$, and is best stated as \emph{training removes excited-state
contamination from the attention field}. The trained attention field is, by that
measure, a better operator than the hand-built smeared basis at every $\beta$
examined, and increasingly so towards the critical point.

What remains open is a volume: both of the largest-$\xi$ ensembles sit at
$\xi \gtrsim L/4$, so the excursion that makes the correlation persuasive is
also the one feature that an $L=32$ run could reassign to finite-volume
squeezing. Until then the excursion is reported, not explained. Beyond that, the
natural extensions are an entropy-matched random control to close the
$\delta A/A$ confound, and the same readout applied to the $SU(2)$ glueball
operator, where the machinery already exists and only the correlation length
does not.

\end{document}
